\documentclass[11pt]{article}
\usepackage[margin=1in]{geometry}
\usepackage[T1]{fontenc}
\usepackage{graphicx}
\usepackage{hyperref}
\usepackage{enumitem}
\setlength{\parskip}{0.5em}
\setlength{\parindent}{0pt}

\title{transformer attention models for language}
\author{arXiv GraphRAG --- phase-6 deep research}
\date{\today}

\begin{document}
\maketitle

\begin{abstract}
The review is excellent and covers the latest theoretical and architectural trends, but misses the critical hardware-aware/IO-aware optimization frontier (e.g., FlashAttention) that defines modern implementation.
\end{abstract}

\section{The Evolution of Attention: From Sequence Foundations to the Redefinition of Transformer Architectures}

The ability to model complex linguistic dependencies has long been the central challenge of natural language processing. While the foundations of sequence modeling and early attention mechanisms provided the initial framework for handling sequential data, the inherent limitations of recurrent structures in capturing long-range context necessitated a fundamental shift in architecture. The importance of mastering attention mechanisms cannot be overstated, as they serve as the primary engine for information retrieval and weighting in modern language models. As these models grow in complexity, understanding the mechanics of how they prioritize information becomes essential for both improving performance and ensuring the reliability of large-scale linguistic agents.

This review traces the historical and technical arc of these developments, beginning with the Transformer revolution and the subsequent rise of transformative pre-training paradigms. We follow this progression into the era of scaling context and implementing structural enhancements designed to handle increasingly massive datasets. The narrative then shifts toward a deeper understanding of model behavior through analyzing and visualizing attention mechanisms, followed by an exploration of hybrid architectures that merge attention with RNNs and convolutions to balance performance and complexity. To address the practical constraints of modern deployment, we examine advancements in efficiency, sparsity, and low-resource optimization. The review concludes by examining the current frontier, where the field is actively deconstructing and redefining attention to move beyond existing architectural constraints.

\subsection{Foundations of Sequence Modeling and Early Attention}

The evolution of sequence modeling transitioned from rigid, fixed-length representations to dynamic, attention-driven architectures. Early efforts focused on utilizing Recurrent Neural Network (RNN) Encoder-Decoder frameworks to learn semantically and syntactically meaningful phrase representations for statistical machine translation [\texttt{arXiv:1406.1078}]. This was further formalized into a general end-to-end approach for sequence-to-sequence learning, where a multilayered Long Short-Term Memory (LSTM) encoder maps input sequences to a fixed-dimensionality vector, which a deep LSTM decoder then processes [\texttt{arXiv:1409.3215}]. While effective, this approach faced a significant architectural bottleneck: the reliance on a single fixed-length vector to represent the entire source sequence [\texttt{arXiv:1409.0473}].

To address this limitation, researchers introduced soft-attention mechanisms that allow the model to automatically search for and align with relevant parts of the source sentence during translation [\texttt{arXiv:1409.0473}]. This shift from a static bottleneck to a joint learning of alignment and translation enabled performance comparable to state-of-the-art phrase-based systems [\texttt{arXiv:1409.0473}]. Subsequent refinements explored the granularity of this mechanism; for instance, distinguishing between global attention, which attends to all source words, and local attention, which focuses on a specific subset [\texttt{arXiv:1508.04025}]. Empirical results demonstrated that local attention could provide a significant 5.0 BLEU point gain over non-attentional systems [\texttt{arXiv:1508.04025}].

The utility of attention extended beyond text-to-text translation into multi-modal and memory-augmented domains. In image caption generation, attention mechanisms were integrated to allow models to "fix their gaze" on salient visual objects during word generation [\texttt{arXiv:1502.03044}]. Simultaneously, the concept of attention was leveraged to build end-to-end memory networks, which utilize multiple computational "hops" over an external memory to reduce the need for supervision in tasks like question answering [\texttt{arXiv:1503.08895}].

Finally, attention mechanisms were optimized for computational efficiency and specific reasoning tasks. In natural language inference (NLI), decomposable attention models were introduced to break problems into parallelizable subproblems [\texttt{arXiv:1606.01933}]. This approach offered a distinct advantage over previous models by using nearly an order of magnitude fewer parameters while remaining trivially parallelizable, even when incorporating intra-sentence attention to capture order information [\texttt{arXiv:1606.01933}]. Collectively, these developments moved the field from simple recurrent mapping toward sophisticated, selective, and efficient information retrieval within sequences.

\subsection{The Transformer Revolution and Pre-training Paradigms}

The landscape of natural language processing was fundamentally altered by the introduction of the Transformer architecture, which replaced traditional recurrent and convolutional layers with a mechanism based solely on self-attention [\texttt{arXiv:1706.03762}]. By dispensing with recurrence, this architecture achieved superior parallelization and significantly reduced training times, setting new benchmarks in machine translation, such as 28.4 BLEU on English-to-German [\texttt{arXiv:1706.03762}]. This shift from sequential processing to attention-based modeling laid the groundwork for the subsequent evolution of deep contextualized word representations.

Early advancements focused on leveraging bidirectional language models (biLM) to capture complex syntax, semantics, and polysemy [\texttt{arXiv:1802.05365}]. This trajectory culminated in BERT, which utilized a bidirectional Transformer encoder pre-trained via Masked Language Model (MLM) and Next Sentence Prediction (NSP) objectives [\texttt{arXiv:1810.04805}]. BERT demonstrated that bidirectional self-attention could replace task-specific architectures through unified fine-tuning, achieving state-of-the-art results on eleven tasks, including an 80.5\% GLUE score [\texttt{arXiv:1810.04805}]. However, BERT's methodology is not without flaws; the use of [MASK] tokens creates a structural mismatch between the pre-training and fine-tuning phases, and larger variants like BERT-large exhibit instability on smaller datasets [\texttt{arXiv:1810.04805}].

Subsequent research sought to address these limitations through optimization and structural refinement. While RoBERTa demonstrated that BERT was significantly undertrained and that hyperparameter optimization and increased data size could yield superior results [\texttt{arXiv:1907.11692}], XLNet attempted to resolve the fundamental pre-train/fine-tune discrepancy and dependency neglect inherent in BERT's masking approach [\texttt{arXiv:1906.08237}]. By employing a permutation language modeling framework, XLNet achieved a generalized autoregressive pre-training that outperformed BERT across 20 tasks [\texttt{arXiv:1906.08237}].

Parallel to these performance gains, researchers explored efficiency and versatility. ALBERT addressed the computational costs of scaling BERT by introducing parameter-reduction techniques, resulting in lower memory consumption and increased training speed while maintaining state-of-the-art performance on benchmarks like GLUE and SQuAD [\texttt{arXiv:1909.11942}]. Diverging from the encoder-only focus, T5 introduced a unified text-to-text framework, converting all language problems into a consistent format to explore the broader landscape of transfer learning [\texttt{arXiv:1910.10683}]. Together, these developments represent a transition from simple sequence transduction [\texttt{arXiv:1706.03762}] to highly optimized, specialized, and unified pre-training paradigms.

\subsection{Scaling Context and Structural Enhancements}

A primary limitation of standard Transformer architectures is the quadratic \$O(n\textasciicircum{}2)\$ scaling of self-attention, which restricts the processing of long sequences due to prohibitive memory and time costs. Research into scaling context has diverged into two main trajectories: enhancing temporal dependency through recurrence and reducing computational complexity through sparse or low-rank attention mechanisms.

To address temporal coherence and context fragmentation, Transformer-XL [\texttt{arXiv:1901.02860}] introduces a segment-level recurrence mechanism and a novel positional encoding scheme. Unlike vanilla Transformers, this approach enables the model to learn long-term dependencies that are 80\% longer than RNNs and 450\% longer than standard Transformers [\texttt{arXiv:1901.02860}]. This architectural modification is highly efficient during evaluation, achieving speeds up to 1,800+ times faster than vanilla Transformers [\texttt{arXiv:1901.02860}].

While Transformer-XL focuses on temporal recurrence, other models target the quadratic complexity of the attention matrix directly. Sparse Transformer [\texttt{arXiv:1904.10509}] proposes sparse factorizations to reduce complexity to \$O(n \textbackslash{}sqrt\{n\})\$ [\texttt{arXiv:1904.10509}]. Building on the necessity for linear scaling, Longformer [\texttt{arXiv:2004.05150}] and BigBird [\texttt{arXiv:2007.14062}] both aim for \$O(n)\$ complexity but employ different structural strategies. Longformer utilizes a combination of local windowed attention and task-motivated global attention [\texttt{arXiv:2004.05150}], whereas BigBird implements a sparse mechanism incorporating \$O(1)\$ global tokens (such as CLS) to maintain universal approximation and Turing completeness [\texttt{arXiv:2007.14062}]. BigBird’s approach allows for handling sequences up to 8x longer than previous methods on similar hardware [\texttt{arXiv:2007.14062}].

Alternative efficiency strategies include low-rank approximations and dynamic projections. Linformer [\texttt{arXiv:2006.04768}] reduces complexity to \$O(n)\$ in both time and space by using a low-rank matrix projection to approximate self-attention [\texttt{arXiv:2006.04768}]. Similarly, the Long-Short Transformer [\texttt{arXiv:2107.02192}] achieves linear complexity by integrating long-range attention with dynamic projection and short-term attention [\texttt{arXiv:2107.02192}]. The Long-Short Transformer demonstrates significant practical utility, handling 3x longer sequences than full-attention on the same hardware and outperforming state-of-the-art models on the Long Range Arena benchmark [\texttt{arXiv:2107.02192}]. While Reformer [\texttt{arXiv:2001.04451}] also seeks to improve efficiency by replacing dot-product attention, the aforementioned models provide diverse structural paths—ranging from the low-rank projections of Linformer [\texttt{arXiv:2006.04768}] to the hybrid local-global attention of Longformer [\texttt{arXiv:2004.05150}] and BigBird [\texttt{arXiv:2007.14062}]—to overcome the fundamental bottlenecks of sequence length.

\subsection{Analyzing and Visualizing Attention Mechanisms}

Research into transformer attention mechanisms has transitioned from investigating the functional necessity of individual heads to developing sophisticated tools for global and word-level interpretability. A central theme in the literature is the inherent redundancy within multi-head attention architectures. Both \texttt{arXiv:1905.09418} and \texttt{arXiv:1905.10650} challenge the assumption that all attention heads are essential for model performance. Specifically, \texttt{arXiv:1905.10650} observes that a large percentage of heads can be removed at test time to improve speed and memory efficiency without significant performance degradation. \texttt{arXiv:1905.09418} advances this by proposing a specific pruning method using stochastic gates and a differentiable relaxation of the L0 penalty. This approach demonstrates that one can remove the vast majority of heads—for instance, 38 out of 48 in an English-Russian neural machine translation task—with a minimal loss of only 0.15 BLEU [\texttt{arXiv:1905.09418}]. While \texttt{arXiv:1905.10650} focuses on the potential for pruning at test time, \texttt{arXiv:1905.09418} provides a more structured architectural modification by identifying specialized, linguistically-interpretable heads that must be preserved [\texttt{arXiv:1905.09418}].

Complementing these efficiency-focused studies is a body of work dedicated to understanding the semantic and syntactic content captured by attention. \texttt{arXiv:1906.04341} provides a foundational analysis of BERT, utilizing an attention-based probing classifier to demonstrate that attention heads capture critical syntactic information, such as coreference, determiners, and direct objects [\texttt{arXiv:1906.04341}]. This analytical depth is further expanded by more recent visualization frameworks. While earlier work often focused on the internal mechanics of the heads themselves, \texttt{arXiv:2305.03210} introduces AttentionViz, an interactive technique that shifts the perspective to a global view. By utilizing joint query-key embeddings, AttentionViz allows researchers to analyze patterns across multiple input sequences in both language and vision transformers [\texttt{arXiv:2305.03210}].

A further refinement in interpretability is seen in the shift from token-level analysis to word-level impact. Whereas many methods examine self-attention scores, the Attention Visualizer package introduced in \texttt{arXiv:2308.14850} focuses on how individual words impact the final embedding representation in encoder-only models [\texttt{arXiv:2308.14850}]. This represents a methodological evolution: where \texttt{arXiv:1906.04341} uses probing to infer linguistic properties [\texttt{arXiv:1906.04341}], and \texttt{arXiv:2305.03210} visualizes the interactions of query and key vectors [\texttt{arXiv:2305.03210}], \texttt{arXiv:2308.14850} seeks to reveal word importance through the lens of the model's final output [\texttt{arXiv:2308.14850}]. Together, these works illustrate a trajectory from proving head redundancy to providing granular, multi-scale tools for explaining transformer decision-making.

\subsection{Hybrid Architectures: Merging Attention with RNNs and Convolutions}

Recent advancements in transformer research have moved toward hybridizing attention mechanisms with recurrent and convolutional structures to address inherent limitations in sequence modeling, such as quadratic complexity, lack of explicit sequentiality, and poor length generalization.

A primary direction involves reintegrating recurrence to capture temporal dependencies. TransfoRNN [\texttt{arXiv:2104.01572}] addresses the fact that standard self-attention lacks explicit token arrangement awareness by cascading recurrent neural networks (RNNs) with self-attention. This architectural modification allows shallow transformer stacks to achieve comparable or superior performance to deeper models, notably reducing parameter counts by up to 10.5\% on Penn Treebank and 27.7\% on WikiText-2 [\texttt{arXiv:2104.01572}]. Similarly, the Recurrent Memory Transformer (RMT) [\texttt{arXiv:2207.06881}] utilizes segment-level recurrent self-attention with specialized memory tokens to manage long-range dependencies. While TransfoRNN focuses on improving the efficiency of shallow models through cascading, RMT optimizes for long-sequence processing, performing on par with Transformer-XL and outperforming it on tasks requiring extended context [\texttt{arXiv:2207.06881}]. Moving beyond simple cascading, Aaren [\texttt{arXiv:2405.13956}] reinterprets the attention mechanism itself as an RNN. By employing a parallel prefix scan algorithm, Aaren enables parallel training while maintaining constant-memory inference, offering significant time and memory efficiency improvements over standard Transformers in low-resource settings [\texttt{arXiv:2405.13956}].

Parallel to recurrent integration, researchers have explored convolutional augmentations to enhance local feature extraction and abstraction. The Convolution-Augmented Transformer (CAT) [\texttt{arXiv:2407.05591}] incorporates convolutional filters directly into the Key (K), Query (Q), and Value (V) embeddings. Unlike the memory-token approach of RMT, CAT leverages convolutional locality to summarize context windows, which allows it to solve associative recall and copying tasks with a single layer and provides guaranteed length generalization [\texttt{arXiv:2407.05591}]. This concept of using convolutions to model structural evolution is further expanded by the "evolving attention" mechanism [\texttt{arXiv:2102.12895}], which integrates residual connections and convolutional layers to allow attention maps across different layers to share knowledge. While CAT uses convolutions to improve specific recall and generalization capabilities [\texttt{arXiv:2407.05591}], the evolving attention mechanism focuses on modeling the evolutionary process of abstraction across hierarchical layers [\texttt{arXiv:2102.12895}].

These hybrid approaches represent a trade-off between global modeling and local efficiency. While the "relaxed attention" method [\texttt{arXiv:2209.09735}] offers a lightweight regularization alternative by smoothing attention weights to prevent overfitting—achieving 26.31\% WER on LRS3—it does not fundamentally alter the architecture's structural relationship with recurrence or convolution. In contrast, the works of TransfoRNN [\texttt{arXiv:2104.01572}], RMT [\texttt{arXiv:2207.06881}], and CAT [\texttt{arXiv:2407.05591}] represent more fundamental shifts toward architectures that synergize the global reach of attention with the inductive biases of RNNs and convolutions.

\subsection{Efficiency, Sparsity, and Low-Resource Optimization}

Recent advancements in transformer architectures have shifted focus from pure scaling toward optimizing computational efficiency, reducing parameter counts, and improving performance in data-constrained environments. A common thread across these works is the mitigation of the quadratic complexity inherent in standard self-attention, though the methodologies vary from structural grouping to hybrid recurrent integration.

To address the high computational requirements of training and dense operations, GroupBERT introduces a decoupled approach by complementing self-attention with a convolutional module and implementing grouped transformations in both feed-forward layers and convolutions [\texttt{arXiv:2106.05822}]. This architectural modification aims to reduce floating-point operations (FLOPs) and time-to-train while maintaining superior performance in language representation learning compared to standard BERT models [\texttt{arXiv:2106.05822}].

While GroupBERT optimizes through structural grouping, other researchers focus on hybridizing transformers with recurrent mechanisms to handle long-sequence inference and low-resource settings. RecurFormer proposes that not all attention heads require full self-attention; instead, it replaces "recency-aware" heads—those focusing on local or short-range dependencies—with linear recurrent neural networks based on the Mamba architecture [\texttt{arXiv:2410.12850}]. This substitution significantly enhances inference efficiency and reduces cache size without the need for token eviction, making it particularly effective for long-input LLM tasks [\texttt{arXiv:2410.12850}]. Similarly, hybrid RNN-Transformer architectures have been employed to optimize language modeling in low-data regimes [\texttt{arXiv:2502.00617}]. By selectively replacing attention layers with feed-forward and quasi-recurrent neural network layers, these models achieve significant reductions in parameter counts and training costs, allowing them to match the performance of much larger models [\texttt{arXiv:2502.00617}].

The challenge of low-resource optimization is also addressed through specialized attention mechanisms designed for specific linguistic constraints. The Attention Link (AL) architecture was developed specifically to improve transformer performance in low training resource scenarios, achieving a state-of-the-art BLEU score of 37.9 on the IWSLT14 de-en task [\texttt{arXiv:2302.00340}]. This contrasts with temporal attention, which optimizes for a different type of "resource": temporal information. By extending the self-attention mechanism to incorporate time points, temporal attention enables models to capture and generalize semantic changes across time-tagged textual data [\texttt{arXiv:2202.02093}].

In summary, the literature reveals a divergence in efficiency strategies: GroupBERT and Conv-Basis [\texttt{arXiv:2405.05219}] target the fundamental quadratic cost of attention, RecurFormer and the hybrid models in [\texttt{arXiv:2502.00617}] leverage recurrence to balance memory and performance, and Attention Link [\texttt{arXiv:2302.00340}] focuses on architectural specialization to overcome data scarcity.

\subsection{The Current Frontier: Deconstructing and Redefining Attention}

Recent research has shifted from merely scaling transformer architectures to fundamentally deconstructing the necessity and efficiency of the attention mechanism. This deconstruction occurs across two dimensions: the functional roles of components and the mathematical principles governing their design.

One line of inquiry investigates whether learnable attention is even essential for sequence modeling. \texttt{arXiv:2506.01115} demonstrates that attention is primarily responsible for in-context reasoning, whereas MLP layers handle knowledge storage. Crucially, they show that even with frozen key and query weights, models can form induction heads and perform competitively, suggesting that the learnable nature of attention may be less critical than its ability to facilitate token mixing. This aligns with findings in \texttt{arXiv:2510.11602}, which identifies token mixing as an indispensable design principle while suggesting that the specific mathematical form and sequence-dependent activations can be relaxed, particularly within hybrid architectures.

As the field moves toward long-context modeling, the focus shifts to overcoming the quadratic complexity of standard self-attention. A comprehensive survey by \texttt{arXiv:2507.19595} categorizes these efforts into linear attention (utilizing kernel approximations or recurrent formulations) and sparse attention (employing fixed patterns or clustering). While this survey provides a rigorous mathematical comparison of update rules and a detailed hardware-centric derivation of chunkwise recurrent representations, it is limited by the omission of certain sections, such as clustering attention, and a lack of comparative performance benchmarks (e.g., perplexity vs. latency).

Contrasting these incremental efficiency gains, \texttt{arXiv:2506.11305} proposes a radical departure from the attention paradigm through the Avey architecture. By replacing self-attention with a ranker paired with an autoregressive neural processor, Avey decouples sequence length from context width, enabling linear scaling and efficient long-range dependency modeling. This stands in opposition to the hybrid approaches explored in \texttt{arXiv:2605.00768}, which seeks to optimize the existing transformer framework by combining global and local attention. Unlike heuristic approaches to local attention, \texttt{arXiv:2605.00768} provides a formal logical characterization, proving that local attention introduces a second temporal operator that strictly enlarges the class of recognizable regular languages beyond what global attention can achieve. This theoretical advantage is empirically supported by a reduction in WikiText-2 perplexity of up to 69.7 relative to global baselines, although the authors acknowledge that their results are limited to lightweight GPT-2-small configurations and may not scale to larger models.

Together, these works suggest that the future of language modeling lies in a transition from monolithic, quadratic attention to specialized, hybrid, or decoupled architectures that prioritize token mixing and expressivity while minimizing computational overhead.

\subsection{Tensions and Open Problems}

A primary tension in the literature concerns the fundamental necessity of the attention mechanism itself for effective sequence modeling. On one hand, the prevailing paradigm relies on the dense, global connectivity of self-attention to capture long-range dependencies, as exemplified by the bidirectional modeling in BERT (\texttt{arXiv:1810.04805}) and the generalized autoregressive approach of XLNet (\texttt{arXiv:1906.08237}). However, recent critical inquiries suggest that the heavy computational tax of attention may be misplaced. For instance, \texttt{arXiv:2506.11305} challenges the indispensability of these mechanisms, arguing that traditional attention-heavy architectures, including the foundational RNN encoder-decoder models (\texttt{arXiv:1406.1078}) and even modern Transformers (\texttt{arXiv:1810.04805}), may be over-engineered for certain linguistic tasks. This debate is further complicated by findings in \texttt{arXiv:2506.01115}, which question whether trainable attention components are strictly necessary, or if random attention patterns could suffice for sequence modeling, thereby suggesting that the "intelligence" of the model may reside in the data or the underlying recurrent structure rather than the attention weights.

A second debate centers on the trade-off between architectural complexity and computational efficiency, particularly regarding how to scale models without incurring quadratic costs. One school of thought focuses on approximating the attention matrix to achieve linear complexity, such as the low-rank projections proposed in Linformer (\texttt{arXiv:2006.04768}). Conversely, other researchers argue that such approximations may sacrifice the very representational power that makes Transformers effective. This tension is visible in the comparison between the efficiency-focused Linformer (\texttt{arXiv:2006.04768}) and the high-capacity BERT (\texttt{arXiv:1810.04805}), as well as in the development of hybrid models. For example, while Longformer (\texttt{arXiv:2004.05150}) seeks to extend context windows through sparse attention, newer research into hybrid RNN-Transformer architectures (\texttt{arXiv:2502.00617}) suggests that the path to efficiency may lie in re-integrating recurrent inductive biases rather than merely sparsifying the attention matrix.

Finally, there is an ongoing disagreement regarding the optimal method for encoding structural and positional information. While early self-attention models relied on absolute or relative position representations to provide order (\texttt{arXiv:1803.02155}), recent work has sought to refine this through more robust masking or structural constraints. For example, Mask Attention Networks (\texttt{arXiv:2103.13597}) critique the limitations of standard relative position representations (\texttt{arXiv:1803.02155}), arguing for more sophisticated attention patterns to strengthen the model's grasp of sequence topology. This struggle to balance positional awareness with computational fluidity remains a core bottleneck in designing architectures that are both structurally expressive and hardware-efficient.

\subsection{Outlook}

The frontier of the field is rapidly shifting away from "attention-at-all-costs" toward a more nuanced synthesis of recurrent and associative memory structures. As we move toward 2026, the focus is transitioning from merely scaling parameters to characterizing the precise expressivity of local versus global attention (\texttt{arXiv:2605.00768}) and deconstructing the design principles that make certain attention patterns effective (\texttt{arXiv:2510.11602}). We anticipate a convergence where the distinction between "Transformer" and "RNN" blurs, leading to hybrid architectures that utilize distilled recurrent structures (\texttt{arXiv:2503.18565}) to achieve the efficiency of linear-time models without losing the deep contextual reasoning that defined the previous era of large language models.

\section*{References}
\begin{enumerate}\small
  \item Learning Phrase Representations using RNN Encoder-Decoder for
  Statistical Machine Translation (2014). \texttt{arXiv:1406.1078}
  \item Neural Machine Translation by Jointly Learning to Align and Translate (2014). \texttt{arXiv:1409.0473}
  \item Sequence to Sequence Learning with Neural Networks (2014). \texttt{arXiv:1409.3215}
  \item Adam: A Method for Stochastic Optimization (2014). \texttt{arXiv:1412.6980}
  \item Neural Turing Machines (2014). \texttt{arXiv:1410.5401}
  \item Neural Machine Translation of Rare Words with Subword Units (2015). \texttt{arXiv:1508.07909}
  \item Deep Residual Learning for Image Recognition (2015). \texttt{arXiv:1512.03385}
  \item Aligning Books and Movies: Towards Story-like Visual Explanations by
  Watching Movies and Reading Books (2015). \texttt{arXiv:1506.06724}
  \item Effective Approaches to Attention-based Neural Machine Translation (2015). \texttt{arXiv:1508.04025}
  \item End-To-End Memory Networks (2015). \texttt{arXiv:1503.08895}
  \item Distilling the Knowledge in a Neural Network (2015). \texttt{arXiv:1503.02531}
  \item Show, Attend and Tell: Neural Image Caption Generation with Visual
  Attention (2015). \texttt{arXiv:1502.03044}
  \item Layer Normalization (2016). \texttt{arXiv:1607.06450}
  \item SQuAD: 100,000+ Questions for Machine Comprehension of Text (2016). \texttt{arXiv:1606.05250}
  \item A Decomposable Attention Model for Natural Language Inference (2016). \texttt{arXiv:1606.01933}
  \item Google's Neural Machine Translation System: Bridging the Gap between
  Human and Machine Translation (2016). \texttt{arXiv:1609.08144}
  \item Pointer Sentinel Mixture Models (2016). \texttt{arXiv:1609.07843}
  \item Gaussian Error Linear Units (GELUs) (2016). \texttt{arXiv:1606.08415}
  \item Using the Output Embedding to Improve Language Models (2016). \texttt{arXiv:1608.05859}
  \item Tying Word Vectors and Word Classifiers: A Loss Framework for Language
  Modeling (2016). \texttt{arXiv:1611.01462}
  \item Attention Is All You Need (2017). \texttt{arXiv:1706.03762}
  \item Convolutional Sequence to Sequence Learning (2017). \texttt{arXiv:1705.03122}
  \item Outrageously Large Neural Networks: The Sparsely-Gated
  Mixture-of-Experts Layer (2017). \texttt{arXiv:1701.06538}
  \item A Broad-Coverage Challenge Corpus for Sentence Understanding through
  Inference (2017). \texttt{arXiv:1704.05426}
  \item BERT: Pre-training of Deep Bidirectional Transformers for Language
  Understanding (2018). \texttt{arXiv:1810.04805}
  \item GLUE: A Multi-Task Benchmark and Analysis Platform for Natural Language
  Understanding (2018). \texttt{arXiv:1804.07461}
  \item Deep contextualized word representations (2018). \texttt{arXiv:1802.05365}
  \item Character-Level Language Modeling with Deeper Self-Attention (2018). \texttt{arXiv:1808.04444}
  \item Self-Attention with Relative Position Representations (2018). \texttt{arXiv:1803.02155}
  \item Scaling Neural Machine Translation (2018). \texttt{arXiv:1806.00187}
  \item Adaptive Input Representations for Neural Language Modeling (2018). \texttt{arXiv:1809.10853}
  \item RoBERTa: A Robustly Optimized BERT Pretraining Approach (2019). \texttt{arXiv:1907.11692}
  \item Transformer-XL: Attentive Language Models Beyond a Fixed-Length Context (2019). \texttt{arXiv:1901.02860}
  \item Generating Long Sequences with Sparse Transformers (2019). \texttt{arXiv:1904.10509}
  \item XLNet: Generalized Autoregressive Pretraining for Language Understanding (2019). \texttt{arXiv:1906.08237}
  \item Exploring the Limits of Transfer Learning with a Unified Text-to-Text
  Transformer (2019). \texttt{arXiv:1910.10683}
  \item ALBERT: A Lite BERT for Self-supervised Learning of Language
  Representations (2019). \texttt{arXiv:1909.11942}
  \item Analyzing Multi-Head Self-Attention: Specialized Heads Do the Heavy
  Lifting, the Rest Can Be Pruned (2019). \texttt{arXiv:1905.09418}
  \item Are Sixteen Heads Really Better than One? (2019). \texttt{arXiv:1905.10650}
  \item What Does BERT Look At? An Analysis of BERT's Attention (2019). \texttt{arXiv:1906.04341}
  \item HuggingFace's Transformers: State-of-the-art Natural Language Processing (2019). \texttt{arXiv:1910.03771}
  \item Transformer++ (2020). \texttt{arXiv:2003.04974}
  \item Learning Hard Retrieval Decoder Attention for Transformers (2020). \texttt{arXiv:2009.14658}
  \item Fixed Encoder Self-Attention Patterns in Transformer-Based Machine
  Translation (2020). \texttt{arXiv:2002.10260}
  \item Memory Transformer (2020). \texttt{arXiv:2006.11527}
  \item Guiding Attention for Self-Supervised Learning with Transformers (2020). \texttt{arXiv:2010.02399}
  \item Language Models are Few-Shot Learners (2020). \texttt{arXiv:2005.14165}
  \item Longformer: The Long-Document Transformer (2020). \texttt{arXiv:2004.05150}
  \item Reformer: The Efficient Transformer (2020). \texttt{arXiv:2001.04451}
  \item An Image is Worth 16x16 Words: Transformers for Image Recognition at
  Scale (2020). \texttt{arXiv:2010.11929}
  \item Transformers are RNNs: Fast Autoregressive Transformers with Linear
  Attention (2020). \texttt{arXiv:2006.16236}
  \item Big Bird: Transformers for Longer Sequences (2020). \texttt{arXiv:2007.14062}
  \item Linformer: Self-Attention with Linear Complexity (2020). \texttt{arXiv:2006.04768}
  \item TransfoRNN: Capturing the Sequential Information in Self-Attention
  Representations for Language Modeling (2021). \texttt{arXiv:2104.01572}
  \item GroupBERT: Enhanced Transformer Architecture with Efficient Grouped
  Structures (2021). \texttt{arXiv:2106.05822}
  \item Evolving Attention with Residual Convolutions (2021). \texttt{arXiv:2102.12895}
  \item Mask Attention Networks: Rethinking and Strengthen Transformer (2021). \texttt{arXiv:2103.13597}
  \item Long-Short Transformer: Efficient Transformers for Language and Vision (2021). \texttt{arXiv:2107.02192}
  \item Relaxed Attention for Transformer Models (2022). \texttt{arXiv:2209.09735}
  \item Multiformer: A Head-Configurable Transformer-Based Model for Direct
  Speech Translation (2022). \texttt{arXiv:2205.07100}
  \item Pre-Training a Graph Recurrent Network for Language Representation (2022). \texttt{arXiv:2209.03834}
  \item Temporal Attention for Language Models (2022). \texttt{arXiv:2202.02093}
  \item Recurrent Memory Transformer (2022). \texttt{arXiv:2207.06881}
  \item Attention Link: An Efficient Attention-Based Low Resource Machine
  Translation Architecture (2023). \texttt{arXiv:2302.00340}
  \item AttentionViz: A Global View of Transformer Attention (2023). \texttt{arXiv:2305.03210}
  \item Stack Attention: Improving the Ability of Transformers to Model
  Hierarchical Patterns (2023). \texttt{arXiv:2310.01749}
  \item Attention Visualizer Package: Revealing Word Importance for Deeper
  Insight into Encoder-Only Transformer Models (2023). \texttt{arXiv:2308.14850}
  \item Knowledge-Infused Self Attention Transformers (2023). \texttt{arXiv:2306.13501}
  \item A Transformer with Stack Attention (2024). \texttt{arXiv:2405.04515}
  \item RecurFormer: Not All Transformer Heads Need Self-Attention (2024). \texttt{arXiv:2410.12850}
  \item On the Power of Convolution Augmented Transformer (2024). \texttt{arXiv:2407.05591}
  \item Attention as an RNN (2024). \texttt{arXiv:2405.13956}
  \item Conv-Basis: A New Paradigm for Efficient Attention Inference and
  Gradient Computation in Transformers (2024). \texttt{arXiv:2405.05219}
  \item Efficient Attention Mechanisms for Large Language Models: A Survey (2025). \texttt{arXiv:2507.19595}
  \item Is Random Attention Sufficient for Sequence Modeling? Disentangling Trainable Components in the Transformer (2025). \texttt{arXiv:2506.01115}
  \item Deconstructing Attention: Investigating Design Principles for Effective Language Modeling (2025). \texttt{arXiv:2510.11602}
  \item Don't Pay Attention (2025). \texttt{arXiv:2506.11305}
  \item Distil-xLSTM: Learning Attention Mechanisms through Recurrent Structures (2025). \texttt{arXiv:2503.18565}
  \item Efficient Language Modeling for Low-Resource Settings with Hybrid
  RNN-Transformer Architectures (2025). \texttt{arXiv:2502.00617}
  \item Characterizing the Expressivity of Local Attention in Transformers (2026). \texttt{arXiv:2605.00768}
\end{enumerate}

\end{document}
